% !TeX root = ../thesis.tex

\chapter{Zusammenfassung}
\label{sec:conclusion_future-work}

\section{Beispiele}
\label{sec:Bsp}


Hier sind einige Beispiele.

\subsection{Abkürzung}
\label{subsec:abkürzung}

So verwende ich eine Abkürzung \gls{acr:djt} und so erneut \gls{acr:djt}.


\subsection{Symbol}
\label{subsec:symbol}
So füge ich ein Symbol ein \gls{symb:h_s}.


\subsection{Bild}
\label{subsec:Bild}
So füge ich ein Bild ein:

\begin{figure}[H]
	\centering
	\includegraphics[width = 1 \textwidth ]{donald.jpg}
	\caption[Titel für das Abbildungsverzeichnis]{Bild Unterschrift \cite{wangVisibleLightCommunications2017}}
	\label{fig:donald}
\end{figure}


\begin{figure}[H]
	\centering
	\includegraphics[width = 1 \textwidth ]{donald.jpg}
	\caption[Titel für das Abbildungsverzeichnis1]{Bild Unterschrift1 \cite{stinyAktiveElektronischeBauelemente2019}}
	\label{fig:donald1}
\end{figure}


So beziehe ich mich auf das Bild ~\ref{fig:donald}

\subsection{Formel}
\label{subsec:formel}

\begin{equation}
	\label{equ:bsp1}
	F_{A} \geq 2 \cdot F_{S}
\end{equation}

So beziehe ich mich auf die Formel ~\ref{equ:bsp1}


\subsection{Tabelle}
\label{subsec:table}


\begin{table}[htb]
	\begin{center}
		\begin{tabular}[h]{c|c|c}	
			
			Stimmen für SleepyJoe & davon gefaked  & Sieger \\
			\hline
			80000000 & 80000000 & Donald \\
			\hline
		\end{tabular}
		\caption{Unterschrift  der Tabelle}
		\label{tab:Tabelle1}
	\end{center}
\end{table}

So beziehe ich mich auf die Tabelle ~\ref{tab:Tabelle1}


\subsection{Referenz und Zitat}
\label{subsec:refcite}
Referenzierung auf ein Kapitel ~\ref{sec:Abschnitt1}. ~\ref{sec:evaluation}
Zitat einfügen

\cite{stinyAktiveElektronischeBauelemente2019}.
\cite{aasdBlabla2021}
So verwende ich eine Online Quelle \gls{online:Eigen}
